% =====================================================================
% EXPERIMENTS
% =====================================================================

\section{Experiments}
\label{sec:experiments}

\begin{figure}[t]
  \centering
  \includegraphics[width=\linewidth]{figures/cifar_results.pdf}
  \caption{CIFAR-10 with 20\% label noise (top) and 100\% label noise
  (bottom), with $K=8$ SGD seeds per width. We plot cross-entropy of the train set, the per-seed
  test set (mean over seeds; band $\pm 1$ s.d.), and the probability-averaged
  ensemble of 8 Gaussian-prior SGLD chains with 10 retained draws per chain
  (green, dashed).
  Under both noise levels the per-seed test cross-entropy exhibits a
  pronounced double-descent peak near the interpolation threshold while the ensemble
  cross-entropy is smaller, particularly near interpolation.
  Test cross-entropy uses independently corrupted labels and 0/1 error uses clean labels.
  We include 0/1 error for reference, but note that it cannot separate near-uniform
  predictions from overconfident wrong predictions in the 100\% noise case---the
  distinction cross-entropy resolves.}
  \label{fig:ensemble}
\end{figure}

We investigate this picture across two model families that span the standard
double-descent literature---CNNs and random forests---each with a varying
parameter count across the interpolation threshold on a classification
task. Throughout we report \emph{cross-entropy} as the
primary metric. Cross-entropy is the natural quantity in our framework because
it estimates the test log-loss $L_n(P^*)$ in~\eqref{eq:testloss} and therefore
measures predictive density rather than 0/1 correctness; classification accuracy
is reported for context but does not enter the decomposition~\eqref{eq:decomp}.

\paragraph{Ensemble as a Bayes predictor.}
Following Section~\ref{sec:gibbs}, we approximate the Bayes predictor
$Q^\pi_{n+1}$ using finite runs of $K$ independent SGLD chains under an explicit Gaussian
prior, averaging retained predictive probabilities within each chain and
then across chains. Crucially, we
average \emph{class probabilities} (the softmax outputs), \emph{not logits}.

\begin{figure}[t]
  \centering
  \includegraphics[width=\linewidth]{figures/gap_vs_seeds.pdf}
  \caption{Effect of the number of seeds $K$ averaged into the SGD ensemble.
  Left panels: the difference in the ensemble test metric between two widths
  ($k{=}8\to 10$ for 20\% noise, $k{=}10\to 14$ for 100\% noise),
  computed as the larger width minus the smaller, as a function of $K$.
  Right panels: the ensemble test curve as $K$ varies. We report 0/1 error
  for 20\% noise and cross-entropy for 100\% noise. Lines give means over
  20 random subsets of 64 seeds (20\% noise) or 16 seeds (100\% noise),
  with no replacement within each subset. Left-panel subsets are paired
  across widths; bands give $\pm 1$ s.d. The full ensemble is used once.
  Under 20\% noise the difference is negative for $K\ge32$.
  Under 100\% noise, the difference remains positive at $K=16$.}
  \label{fig:gap-vs-seeds}
\end{figure}

\subsection{Bayes predictor and double descent}
\label{sec:exp-ensemble}

We train a 5-layer CNN of varying width $k$ on CIFAR-10 with a fixed
fraction of training labels replaced by uniformly random labels, sweeping widths. For each
$k$ we train $K=8$ independently initialized SGD seeds and $K=8$ SGLD
chains, and compare the mean per-seed SGD test cross-entropy with the
SGLD ensemble test cross-entropy. SGD is a finite optimization of the
maximum-likelihood objective. Full
training details (architecture, optimizer, schedule, batch size, number
of gradient steps) are deferred to Appendix~\ref{app:training}.

\paragraph{CIFAR-10.}
We focus first on the canonical $20\%$ label noise setup
(Figure~\ref{fig:ensemble}). The training cross-entropy approaches
zero at the interpolation threshold and remains close to zero for all larger $k$,
confirming that we are operating in the interpolation regime. The per-seed
test cross-entropy exhibits the double-descent peak near the interpolation
threshold, while the SGLD ensemble cross-entropy
shows no pronounced interpolation peak. We hypothesize that
the second descent is primarily produced by the MLE-vs-Bayes gap~(III),
but cannot conclude this directly due to the possible confounding effect
of the approximation-and-shift term~(I).

\paragraph{Random labels.}
To isolate the MLE-vs-Bayes gap~(III) from the approximation-and-shift
term~(I), we repeat the experiment with $100\%$ label noise: training labels
drawn uniformly at random, independently of the input. Under this $P^*$ the
true conditional is uniform, $P^*(Y\mid X)\equiv 1/10$. Setting the final-layer weights and biases to zero gives uniform
softmax outputs, so there is a $\theta_0\in\Theta$ with
$D(P^*\,\|\,P_{\theta_0})=0$. Hence
$P^*_\Theta=P^*$ in distribution and
$(\mathrm{I})=\mathcal{R}_n(P^*)-\mathcal{R}_n(P^*_\Theta)=0$.

The training cross-entropy approaches zero at the interpolation threshold, and
the per-seed test cross-entropy exhibits a pronounced peak there
(Figure~\ref{fig:ensemble}). The ensemble test cross-entropy remains near
$\log 10$, so the empirical SGD-vs-ensemble gap is large near the
interpolation threshold. At $k=14$, the mean SGD loss is $8.659$ nats and
the SGLD ensemble loss is $2.303$ nats, giving a gap of $6.356$ nats.
Since $P^*_\Theta=P^*$ here by construction,
\eqref{eq:gap-lower-bounds-III} identifies the expected MLE-vs-ensemble gap
as a lower bound on (III) in the ideal posterior-sampling model.
The large measured gap near interpolation supports the MLE-vs-Bayes gap
as a major contributor to the peak.

Longer chains and a smaller step size changed ensemble log-loss by at most
$0.019$ nats at the checked widths (Appendix~\ref{app:sampling}).
These checks preserve the qualitative comparison. An SGD-ensemble
control is also included there.

\paragraph{SGD ensembling.}
As a practical extension, we also average the class probabilities of
independently trained SGD models.
Figure~\ref{fig:gap-vs-seeds} compares the ensemble curves as $K$ varies.
Increasing $K$ improves the ensemble test metric. Under $20\%$ noise,
the plotted ensemble curves are monotone for $K=32$ and $K=64$; under
$100\%$ noise, the peak shrinks but remains at $K=16$.

\begin{figure}[t]
  \centering
  \includegraphics[width=\linewidth]{figures/rf_double_descent.pdf}
  \caption{Random forests and BART on MNIST ($n=10^4$, 10 classes)
  across the two-phase complexity axis $N^{\mathrm{max}}_{\mathrm{leaf}}/N_{\mathrm{tree}}$.
  \textbf{Top row}: clean training labels. \textbf{Bottom row}: $100\%$ random
  training labels. Both rows use clean test labels. Solid curves give the
  original single-forest train and test scores. The BART ensemble (green, dashed)
  averages regression outputs over 16 chains with 25 draws each, then clips
  and normalizes the ten scores for evaluation. Under both label settings the
  forest test cross-entropy (orange) exhibits a pronounced double-descent peak
  near the interpolation threshold. With random training labels, the BART
  ensemble cross-entropy stays near $\log 10$.}
  \label{fig:rf}
\end{figure}

\subsection{Random forests and BART}
\label{sec:exp-rf}

To compare with a structurally different model family, we now
consider random forests \citep{breiman2001random} on MNIST, following the
double-descent setup of \citet{belkin2019reconciling}. We train on
$n=10^4$ samples and sweep complexity along the two-phase axis
$N^{\mathrm{max}}_{\mathrm{leaf}}/N_{\mathrm{tree}}$: Phase~1 holds
$N_{\mathrm{tree}}=1$ and varies the maximum number of leaves
$N^{\mathrm{max}}_{\mathrm{leaf}}$ from small values up to the
interpolation threshold (where the train cross-entropy first reaches zero);
Phase~2 fixes $N^{\mathrm{max}}_{\mathrm{leaf}}$ at the interpolation
threshold and grows $N_{\mathrm{tree}}$. We run the experiment under both
clean labels and $100\%$ random labels. For comparison, we use a BART-style Gaussian regression ensemble
\citep{chipman2010bart}, averaging 16 chains with 25 retained draws each
before clipping and normalizing the ten class scores. This is an empirical
comparison, not a classification posterior predictive to which
\eqref{eq:gap-lower-bounds-III} directly applies.

Figure~\ref{fig:rf} shows train and test cross-entropy across the
complexity axis under both label settings. The forest test cross-entropy
exhibits a pronounced double-descent peak around the interpolation
threshold and then drops sharply once additional trees are added in
Phase~2; the peak is more dramatic under $100\%$ random labels. The
BART ensemble cross-entropy is substantially smaller near interpolation,
mirroring the attenuation of the peak in the CNN result.

